\chapter{Reumatoidná artritída}

1 Reumatoidná artritída Reumatoidná artritída (RA) predstavuje závažné, chronické zápalové ochorenie au- toimunitnej etiológie, ktoré sa vyznacuje systémovým charakterom a progresívnym priebe- hom. Jablkom sváru v patogenéze tohto ochorenia je imunitný systém, ktorý omylom iden- tifikuje vlastné tkanivá ako cudzorodé, co vedie k deštrukcii, bolesti a strate funkcie. Hoci primárnymciel’ompatologickéhoprocesusúperiférnesynoviálneklby,najmädrobnésklbe- nia rúk a nôh, RA sa v plnej miere prejavuje ako systémová choroba postihujúca aj mimo- klbové štruktúry a vnútorné orgány. Tento aspekt ochorenia významne prispieva k zvýšenej morbiditeaskráteniudlžkyživotapostihnutýchjedincov.Ochoreniesamôžemanifestovat’v ktoromkol’vekveku,avšakdemografickékrivkyjasnepoukazujúnavrcholincidencievpro- duktívnomveku,typickymedzi30.aža50.rokomživota,cománesmiernycelospolocenský a ekonomický dopad. Etiopatogenéza ochorenia je multifaktoriálna a hoci nie je do posled- ného detailu objasnená, predpokladá sa zložitá interakcia genetickej predispozície hostitel’a a spúšt’acích faktorov vonkajšieho prostredia. Priebeh choroby je vysoko individuálny -- od miernych foriem s pomalou progresiou až po agresívne formy vedúce k rýchlej invalidizácii (Rovenský,2001). Anatomickým substrátom, kde sa odohráva hlavná dráma autoimunitnej reakcie, je synoviálna membrána. Za normálnych okolností tenká vrstva buniek sa pri RA mení na hy- perplastické,zapálenétkanivoprestúpenébunkamiimunitnéhosystému.Tátozmenenásyno- via, oznacovaná ako panus, sa správa ako lokálny tumor, ktorý invazívne prerastá do okolia, nahlodávaklbovúchrupavkuanásledneajsubchondrálnukost’.Výsledkomtohtoprocesuje nielen bolest’ a opuch, ale predovšetkým nezvratná deštrukcia klbovej architektúry a strata hybnosti(Rovenský,2001;Shahetal.,2024).


\section{Funkcnáanatómiaklbovéhosystému}

Pochopenie anatomických a biomechanických princípov klbového aparátu je nevy- hnutným predpokladom pre správnu interpretáciu patologických zmien pri reumatoidnej ar- tritíde.Klbyl’udskéhotelamožnorozdelit’podl’arôznychkritérií,pricomzpohl’adureuma- tológie má primárny význam delenie podl’a prítomnosti klbovej dutiny na spojenia nepravé (synartrózy)apravéklby(diartrózy).Právepravéklbyobsahujúcesynoviálnumembránusú tercomautoimunitnéhoútokupriRA(Dylevský,2009).


\subsection{Stavbasynoviálnehoklbu}

Synoviálny klb predstavuje zložitú anatomickú jednotku, ktorej základnými kompo- nentmi sú klbové povrchy pokryté hyalínnou chrupavkou, klbové puzdro s jeho vrstvami a klbová dutina vyplnená synoviálnou tekutinou. Každá z týchto štruktúr zohráva špecifickú úlohu v biomechanike pohybu a súcasne predstavuje potenciálny substrát pre patologické 2 zmeny(Dylevský,2009). Klbová chrupavka pokrýva artikulacné plochy kostí a zabezpecuje hladký pohyb s minimálnymtrením.JetvorenáprevažneextracelulárnoumatrixbohatounakolagéntypuIIa proteoglykány,ktoréviažuvoduaposkytujúchrupavkejejtypickéviskoelastickévlastnosti. Chondrocyty, jediné bunky prítomné v chrupavke, udržiavajú rovnováhu medzi syntézou a degradácioumatrix.PriRAdochádzaknarušeniutejtorovnováhyvprospechdegradácie,co vediekprogresívnemuúbytkuchrupavkovéhotkaniva. Klbovépuzdropozostávazdvochvrstiev.Vonkajšiafibróznavrstva(membranafib- rosa) je tvorená hustým väzivom bohatým na kolagénové vlákna typu I, ktoré zabezpecujú mechanickú stabilitu klbu. Vnútorná synoviálna vrstva (membrana synovialis) je pri RA miestom najintenzívnejších zápalových zmien. Za fyziologických podmienok je synoviálna membránatenká,tvorenájednouažtromivrstvamibuniek--synoviocytov(Kolár,2009).


\subsection{Synoviálnamembránaajejúloha}

Synoviálna membrána je metabolicky aktívne tkanivo, ktoré produkuje synoviálnu tekutinu -- viskózny ultrafiltrátt krvnej plazmy obohatený o kyselinu hyalurónovú. Táto te- kutinaplníniekol’kokl’úcovýchfunkcií:lubrikáciuklbovýchpovrchov,výživuavaskulárnej chrupavkydifúziouživínaodstranovaniemetabolickýchproduktovzklbovejdutiny. V synoviálnej membráne rozlišujeme dva hlavné typy buniek. Synoviocyty typu A majúmakrofágovýpôvodazabezpecujúfagocytózudebrisaimunitnýdohl’ad.Synoviocyty typu B sú fibroblastového charakteru a produkujú komponenty synoviálnej tekutiny, najmä kyselinu hyalurónovú a lubricín. Pri RA dochádza k masívnej infiltrácii synoviálnej mem- brányimunitnýmibunkamiatransformáciirezidentnýchsynoviocytovdoagresívnehofeno- typu(Smith,2015).


\subsection{AnatomicképredilekcnémiestapriRA}

Reumatoidná artritída vykazuje charakteristickú predilekciu pre urcité klbové sku- piny, co má diagnostický význam. Najcastejšie postihnuté sú metakarpofalangeálne (MCP) klby,proximálneinterfalangeálne(PIP)klbyrúk,zápästnéklbyametatarzofalangeálne(MTP) klby nôh. Distálne interfalangeálne (DIP) klby sú typicky ušetrené, co odlišuje RA od oste- oartrózyapsoriatickejartritídy(Olejárová,2008). Táto selektívna distribúcia postihnutia nie je náhodná. Drobné klby prstov obsahujú pomerne vel’kú synoviálnu membránu v pomere k vel’kosti klbu, co poskytuje väcšiu plo- chupreautoimunitnúreakciu. Navyšetietoklbysúvystavenéopakovanému mechanickému stresu pri bežných denných cinnostiach, co môže potencovat’ zápalový proces (K. Pavelka; Rovenskýetal.,2003). 3


\section{Epidemiológia}

Výskyt a rozšírenie RA sa líši podl’a diagnostických schém, pricom incidencia je vyššia u žien než u mužov, a to v 70 percentách prípadov. Prevalencia sa pohybuje okolo pol percenta až jedného percenta populácie, líši sa avšak podl’a regiónu a etnických skupín. Zatial’ co v Amerike je prevalencia vyššia, v Afrike je práve naopak nižšia. RA je charakte- ristickávýskytomvoveku20až55rokov;snovýmištúdiamisapotvrdilo,ževznikáužajv neskoršomštádiuveku(Rovenský,2001;Rovenský,1998). Podl’a odhadov v roku 2020 trpelo RA celosvetovo približne 17,6 milióna obyvate- l’ov, po zohl’adnení vekovej štruktúry populácie sa odhaduje, že celosvetovo postihuje re- umatoidnáartritídapribližne0,21%l’udí.Vporovnanísrokom1990došlok14,1%nárastu poctuprípadovRA.Predpokladása,žedo25rokovpocetpacientovsRAstúpnena31,7mi- lióna. Epidemiologické štúdie v Cíne uvádzajú okolo 5 miliónov pacientov, co predstavuje prevalenciu0,42%populácie(Gaoetal.,2020;Rovenský,2001;Rovenský,1998). U starších pacientov zaznamenávame odlišné epidemiologické charakteristiky. Ge- rontorevmatológia poukazuje na špecifiká RA v geriatrickej populácii, kde ochorenie casto nadobúda atypický klinický priebeh. Séropozitivita býva nižšia a diferenciálna diagnostika vocipolymyalgiireumatikeaosteoartrózejenárocnejšia(Rovenský,2014).


\section{Etiológiaapatofyziológia}

Súcasná veda nazerá na vznik reumatoidnej artritídy ako na viacstupnový proces, v ktorom dochádza k prelomeniu imunologickej tolerancie. Na zaciatku stojí geneticky pre- disponovaný jedinec (nosic urcitých alel HLA-DRB1), u ktorého environmentálny spúš- t’ac --napríkladfajcenie,chronickáparodontitída(baktériaPorphyromonasgingivalis)alebo zmenyvcrevnommikrobióme--naštartujeprodukciuautoprotilátok.Tentopredklinickýstav môžetrvat’rokypredtým,nežsaobjavíprvásynovitída.Kl’úcovýmibiomarkerovýmihrácmi v tomto procese sú reumatoidný faktor (RF) a protilátky proti cyklickému citrulínovanému peptidu (anti-CCP), ktorých prítomnost’ casto predchádza klinickým symptómom (Shah et al.,2024;Schereretal.,2020).


\subsection{Autoimunitnémechanizmy}

Autoimunita predstavuje stav, pri ktorom imunitný systém stráca schopnost’ rozli- šovat’ medzi vlastnými a cudzorodými antigénmi. Pri fyziologických podmienkach existujú viaceré kontrolné mechanizmy -- centrálna a periférna tolerancia -- zabranujúce útoku na vlastné tkanivá. U pacientov s RA dochádza k zlyhaniu týchto mechanizmov na viacerých úrovniach(Buc,2016). Genetická predispozícia zohráva kl’úcovú úlohu v rozvoji autoimunitného procesu. Asociácia s hlavným histokompatibilným komplexom (MHC), konkrétne s urcitými ale- lami HLA-DRB1 obsahujúcimi tzv. zdiel’aný epitop (shared epitope), je jedným z najsil- nejších genetických rizikových faktorov. Tieto molekuly HLA prezentujú autoantigény T- 4 lymfocytom,címiniciujúpatologickúimunitnúodpoved’(Taketal.,2011). Proces citrulinácie predstavuje ústredný mechanizmus v patogenéze RA. Ide o po- sttranslacnú modifikáciu proteínov, pri ktorej enzým peptidylarginín deimináza (PAD) kon- vertuje arginín na citrulín. Citrulínované proteíny sa stávajú neoantigénmi, voci ktorým or- ganizmus vytvára protilátky (ACPA). Tieto autoprotilátky sa môžu objavit’ roky pred mani- festáciouklinickýchpríznakovaichprítomnost’jesilnýmprediktoromerozívnehopriebehu ochorenia(Choy,2012).


\subsection{Patofyziologickákaskáda}

V patofyziologickej kaskáde dominuje dysregulácia imunitného systému. Dochádza k patologickej aktivácii T-lymfocytov (najmä Th1 a Th17 subpopulácií), ktoré migrujú do klbovej výstelky. Tu prostredníctvom medzibunkovej komunikácie stimulujú makrofágy, B- lymfocyty a fibroblasty k nadprodukcii prozápalových mediátorov. Vzniká tzv. cytokínová búrka, v ktorej dominujú tumor nekrotizujúci faktor alfa (TNF-α), interleukín-1 (IL-1) a interleukín-6 (IL-6). Tieto molekuly sú zodpovedné nielen za lokálny zápal a bolest’, ale aj za systémové prejavy, ako je únava, horúcka ci tvorba reaktantov akútnej fázy v peceni (CRP)(Gaoetal.,2020).


\subsection{Úlohasynoviálnychfibroblastov}

Synoviálne fibroblasty,oznacované ajako fibroblastom podobné synoviocyty(FLS), zohrávajúcentrálnuúlohuvpatogenézeRA.Narozdielodimunitnýchbuniek,ktorédoklbu infiltrujú zvonka, FLS sú rezidentné bunky synoviálnej membrány, ktoré pri RA podliehajú fenotypovej transformácii. Nadobúdajú agresívne, nádoru podobné vlastnosti charakterizo- vanézvýšenouproliferáciou,zníženouapoptózouaschopnost’ouinvadovat’dochrupavkya kosti(Smith,2015). Aktivované FLS produkujú široké spektrum deštruktívnych enzýmov, predovšetkým matrix metaloproteinázy (MMP-1, MMP-3, MMP-13), ktoré degradujú kolagén a d’alšie komponenty extracelulárnej matrix. Súcasne sekretujú chemokíny prit’ahujúce d’alšie zápa- lovébunkyarastovéfaktorypodporujúceangiogenézuvzapálenomtkanive.Tentoautokriný aparakrinýkomunikacnýsystémvytvárasebaposilnujúcicykluschronickéhozápalu(Smith, 2015;Taketal.,2011).


\subsection{Tvorbapanusuadeštrukciaklbu}

Chronická stimulácia vedie k transformácii synoviocytov, ktoré sa zacínajú nekon- trolovane množit’ a správat’ invazívne. Vytvára sa hypervaskularizované granulacné tka- nivo -- panus. Bunky panusu produkujú metaloproteinázy a iné enzýmy, ktoré doslova "trá- via"klbovú chrupavku. Súcasne dochádza k aktivácii osteoklastov, buniek odbúravajúcich kost’, co vedie k vzniku typických kostných erózií. Tento deštruktívny proces je nezvratný; raz znicenú chrupavku a kost’ organizmus nedokáže plnohodnotne regenerovat’,co podciar- 5 kujenutnost’vcasnejaagresívnejliecby(Rovenský,2001;Isaacsetal.,2002).


\section{Klinickýobraz}

Klinická manifestácia RA je povestná svojou heterogenitou. U väcšiny pacientov (približne 70\%) sa ochorenie vkráda do života plazivo, nenápadne, v priebehu niekol’kých týždnovcimesiacov.Prvýmipríznakminemusiabyt’bolestiklbov,alenešpecifické,prodro- málne symptómy pripomínajúce virózu: nmerná únava, subfebrílie, nechutenstvo a mierny úbytok hmotnosti. Až následne sa pridružuje typická artikulárna symptomatológia. Domi- nantným znakom je ranná stuhnutost’ klbov. Na rozdiel od krátkodobej stuhnutosti pri ar- tróze, pri RA tento stav trvá dlhšie než hodinu, casto celé predpoludnie, a odznieva až po rozhýbaní. Stuhnutost’ je odrazom nahromadenia zápalového edému v klbe pocas nocného pokoja(Nemec,2023). Artritídasaklinickyprejavujeklasickoupäticoupríznakovzápalu:bolest’ou(dolor), opuchom (tumor), zvýšenou teplotou klbu (calor) a poruchou funkcie (functio laesa); zacer- venanie (rubor) býva pri RA menej casté. Postihnutie je typicky polyartikulárne (zasahuje viac ako 3 klbové oblasti) a symetrické (postihuje rovnaké klby na oboch stranách tela). Predilekcnýmilokalitami súdrobnéklbyruky --metakarpofalangeálne(MCP) aproximálne interfalangeálne (PIP) klby, ako aj zápästia a metatarzofalangeálne (MTP) klby nohy. Dis- tálne interfalangeálne klby (DIP) bývajú typicky ušetrené, co je dôležitý diferenciacný znak oprotiosteoartrózeapsoriatickejartritíde(Nemec,2023;Rovenský,2001). U menšej casti pacientov môže mat’ choroba akútny, búrlivý nástup, ktorý pacienta doslova pripúta na lôžko v priebehu niekol’kých dní. Ešte zriedkavejší je palindromatický reumatizmus, charakterizovaný krátkymi, opakujúcimi sa záchvatmi artritídy, ktoré nezane- chávajú trvalé následky, no môžu prejst’ do klasickej RA. Zvláštnu pozornost’ si vyžaduje RA u starších l’udí, ktorá sa casto prezentuje asymetrickým postihnutím vel’kých klbov, najmäramennýchpletencov(tzv."shoulder-likeönset),comôževiest’kzámenespolymyal- gioureumaticou(Nemec,2023).


\subsection{Klbovédeformityaštrukturálnezmeny}

Dlhodobo neostatocne kontrolovaný zápal vedie k deštrukcii väzivového aparátu a klbových puzdier, co má za následok vznik typických deformít. Tieto zmeny nie sú len koz- metickým defektom, ale vedú k t’ažkej funkcnej poruche úchopu a lokomocnej schopnosti (Olejárová;kol.,2016). Rukaazápästie:ArtikulárnykomplexrukyjezrkadlomRA.Synovitídazápästiave- diekdeštrukciikarpálnychkosticiekaväzov,cocastovyústidovolárnejsubluxáciezápästia a instability. Castou komplikáciou je kompresia *nervus medianus* -- syndróm karpálneho tunela. Na prstoch pozorujeme charakteristickú ulnárnu deviáciu (vybocenie prstov smerom klakt’ovejkosti)vMCPklboch.Dalšímitypickýmideformitamisú: • Deformitagombíkovejdierky(Boutonnière):Jecharakterizovanáfixovanouflexiou 6 vPIPklbeahyperextenziouvDIPklbe.Vznikápretrhnutímcentrálnehopruhuexten- zorovejšl’achy. • Deformita labutej šije (Swan neck): Opacná situácia -- hyperextenzia v PIP klbe a flexia v DIP klbe, spôsobená spazmom vnútorných svalov ruky a laxicitou volárnej platnicky. Noha: Postihnutie nôh je casté a bolestivé. Dochádza k poklesu priecnej klenby, vznikuhalluxvalgus(vybocenýpalec)akladivkovýchprstov(digitimallei).Podhlavickami metatarzov sa tvoria bolestivé otlaky, ktoré pacientovi st’ažujú chôdzu a vyžadujú ortope- dickúobuv. Vel’ké nosné klby: Koleno je postihnuté casto, s tvorbou rozsiahlych výpotkov. Hy- pertrofickásynoviaalebotekutinasamôžepretlacit’dopodkolennejjamkyavytvorit’Bake- rovucystu,ktorejprasknutieimitujetrombózužíllýtka.Bedrovýklbbývapostihnutýmenej casto,alekoxitídamátendenciukrýchlejprogresiiavyžadujecastototálnuendoprotézu. Chrbtica:Kritickouoblast’oujekrcnáchrbtica.Zápalvokolíatlantoaxiálnehosklbe- nia(C1-C2)môževiest’kuvol’neniuväzovainstabilite.Tentostavježivotohrozujúci,pre- tože pri neopatrnom pohybe môže dôjst’ k útlaku miechy. Každý pacient s RA podstupujúci operáciuvcelkovejanestéziimusímat’RTGkrcnejchrbticenavylúcenietejtokomplikácie. Obr. 1.1: Schematické znázornenie typických deformít ruky pri pokrocilej reumatoidnej ar- tritíde.ZobrazenájeulnárnadeviáciavMCPklbochadeformityprstov(Netter,2014) 7


\subsection{Mimoklbové(extraartikulárne)prejavy}

Reumatoidnáartritídajesystémovéochorenieparexcellence.Prítomnost’mimoklbo- výchprejavovsignalizujet’ažšiuformuochoreniaahoršiuprognózu(Rovenský,2001). Reumatoidné uzly: Sú najcastejším extraartikulárnym prejavom, vyskytujúcim sa u 20-30\% séropozitívnych pacientov. Sú to tuhé, nebolestivé podkožné útvary, ktoré sa tvoria najmä v miestach mechanického tlaku -- nad lakt’ovým výbežkom (olecranon), na chrbte rukycinaachilovke. Ocné postihnutie: Najbežnejšia je keratokonjunktivitída sicca (syndróm suchého oka) ako súcast’ sekundárneho Sjögrenovho syndrómu. Závažnejšou komplikáciou je sk- leritídaaleboepiskleritída,ktorásaprejavujebolestivýmzacervenanímokaamôževiest’až kpoškodeniuzraku. Pl’úcne postihnutie: Pl’úca sú castým tercom. Môže íst’ o pleuritídu (zápal pohrud- nice) s výpotkom, alebo o závažnejšiu intersticiálnu pl’úcnu fibrózu, ktorá postupne znižuje dýchaciukapacitupl’úc. Vaskulitída: Reumatoidná vaskulitída je zápal ciev, ktorý môže viest’ k tvorbe kož- ných vredov, nekróz na prstoch alebo k postihnutiu nervov (mononeuritis multiplex). Je to vážnakomplikáciavyžadujúcaagresívnuliecbu.


\section{Diagnostika}

Modernáreumatológiakladiedôraznavcasnúdiagnostiku,tzv.öknopríležitosti"(window ofopportunity).Ciel’omjezachytit’ochorenievjehopociatkoch,kedyjeterapeutickáinter- vencia najúcinnejšia a dokáže zabránit’ trvalému poškodeniu klbov. Diagnostický proces je komplexnýavychádzazosyntézyklinickéhoobrazu,laboratórnychvýsledkovazobrazova- cíchmetód.ZlatýmštandardomsúvsúcasnostiklasifikacnékritériáACR/EULARzroku 2010,ktorébolivyvinutépráveprezáchytvcasnýchštádiíartritídy(Aletahaetal.,2010). Bodovací systém kritérií hodnotí štyri domény, pricom na potvrdenie diagnózy "jed- noznacnejRA"jepotrebnýsúcet6aviacbodovzmaximálnych10: 1. Postihnutie klbov (0 -- 5 bodov): Dôraz sa kladie na pocet a typ postihnutých klbov. Najvyššie bodové ohodnotenie (5 bodov) získava pacient s postihnutím viac ako 10 klbov,pricomaspon jedenznichmusíbyt’malýklb. 2. Sérológia (0 -- 3 body): Prítomnost’ reumatoidných faktorov (RF) a protilátok proti citrulínovaným peptidom (ACPA/anti-CCP). Vysoké titre týchto protilátok sú vysoko špecificképreRAaprognostickynepriaznivé. 3. Reaktanty akútnej fázy (0 -- 1 bod): Laboratórny dôkaz systémového zápalu -- zvý- šenásedimentáciaerytrocytov(FW)aleboC-reaktívnyproteín(CRP). 4. Trvaniesymptómov(0--1bod):Pretrvávaniet’ažkostídlhšieako6týždnovodlišuje chronickúartritíduodprechodnýchvírusovýchartralgií. 8 Okrem klasického röntgenového vyšetrenia (RTG), ktoré v skorom štádiu môže byt’ negatívne,sadnesdopoprediadostávajúsenzitívnejšiemetódy.Ultrasonografia(USG)klbov s Power Doppler módom dokáže zobrazit’ synoviálny zápal a prekrvenie skôr, než dôjde k poškodeniukosti.Magnetickárezonancia(MRI)jenajcitlivejšoumetódounadetekciukost- néhoedému,ktorýjepredzvest’oubudúcicherózií(Smolenetal.,2018).


\section{Komplikáciereumatoidnejartritídy}

Jenevyhnutnévnímat’RAnielenakoizolovanúartropatiu,aleakokomplexnúsysté- movú chorobu, ktorej zápalový potenciál indukuje rad závažných komorbidít a komplikácií, majúcich zásadný vplyv na morbiditu a predcasnú mortalitu postihnutých jedincov. Tieto nežiaduce stavy môžu pramenit’ priamo z aktivity základného ochorenia alebo môžu byt’ iatrogénnymnásledkomdlhodobejimunosupresívnejliecby(Smolenetal.,2018).


\subsection{Kardiovaskulárnekomplikácie}

Epidemiologické štúdie konzistentne potvrdzujú, že populácia s RA celí markantne vyššiemurizikukardiovaskulárnychpríhod(približne1,5až2-násobne)vporovnanísozdra- vou populáciou. Práve kardiovaskulárne zlyhanie je hlavnou prícinou skrátenia dlžky života u týchto pacientov. Pretrvávajúci systémový zápal poškodzuje endotelovú výstelku ciev a akceleruje aterosklerotické procesy, co vyúst’uje do zvýšenej incidencie infarktu myokardu, cievnychmozgovýchpríhodakongestívnehosrdcovéhozlyhania(Smolenetal.,2018). Mechanizmykardiovaskulárnehopoškodenia Patofyziologický vzt’ah medzi reumatoidnou artritídou a kardiovaskulárnym ocho- rením je komplexný a zahrna viaceré mechanizmy. Chronický systémový zápal prispieva k endotelovejdysfunkcii--pociatocnémuštádiuaterosklerózy.Prozápalovécytokíny(TNF-α, IL-6) aktivujú endotelové bunky, cím podporujú adhéziu monocytov a ich transmigráciu do cievnej steny. Výsledkom je formácia aterómových plátov, ktoré sú u pacientov s RA casto menejstabilnéanáchylnejšienaruptúru(Freedmanetal.,2014). Rizikový profil je navyše zhoršený prítomnost’ou konvencných faktorov, ako sú hy- pertenziaadyslipidémia,ktorémôžubyt’potencovanéterapioukortikoidmianesteroidnými antireumatikami(NSA).Zklinickéhohl’adiskajevýznamnáajperikarditída,ktoráhocicasto prebieha asymptomaticky, môže byt’ detegovaná echokardiograficky. Striktný manažment krvnéhotlakualipidovjepretoimperatívom(Smolenetal.,2018;Freedmanetal.,2014). Preventívnestratégie PrevenciakardiovaskulárnychpríhodupacientovsRAvyžadujekomplexnýprístup. Na prvom mieste je efektívna kontrola zápalovej aktivity základného ochorenia -- pacienti v remisii majú významne nižšie kardiovaskulárne riziko. Dalej je potrebný striktný manaž- menttradicnýchrizikovýchfaktorovvrátaneliecbyhypertenzie,dyslipidémieadiabetu.Od- 9 porúca sa používat’ modifikované skórovacie systémy (SCORE) s korekcným faktorom 1,5 prepacientovsRA.Významnájeajúlohapravidelnejfyzickejaktivity,ktoránielenzlepšuje kardiovaskulárnukondíciu,alemáajprotizápalovéúcinky(Freedmanetal.,2014).


\subsection{Osteoporózaafraktúry}

Demineralizáciakostnéhotkaniva,osteoporóza,predstavujejednuznajcastejšíchex- traartikulárnych komplikácií s komplexnou etiológiou. V patogenéze rozlišujeme lokálnu periartikulárnu osteoporózu, ktorá je priamym následkom pôsobenia prozápalových cyto- kínov v okolí postihnutého klbu, a generalizovanú systémovú osteoporózu. Na jej vzniku sa podiel’a synergický efekt chronického zápalu, imobilizácie pacienta (zníženie mechanic- kej zát’aže kosti pre bolest’) a najmä dlhodobej kortikoterapie, ktorá inhibuje osteoblastickú aktivitu a redukuje intestinálnu resorpciu vápnika. Dôsledkom je zvýšená fragilita skeletu a vysokérizikopatologickýchfraktúr,predovšetkýmstavcovaproximálnehofemuru,covedie kd’alšiemuzhoršeniufunkcnéhostavu(KarelPavelkaetal.,2018).


\subsection{Infekcnékomplikácie}

Dysfunkcia imunitného dohl’adu, ktorá je inherentnou súcast’ou patofyziológie RA, v kombinácii s úcinkami potentných imunosupresív (vrátane biologík), robí pacientov zra- nitel’nými voci infekcným agens. Infekcie respiracného traktu (pneumónie), urogenitálneho systému a kože majú u týchto pacientov tendenciu k t’ažšiemu priebehu a castejším kom- plikáciám. Pri liecbe inhibítormi TNF-α existuje špecifické riziko reaktivácie latentnej tu- berkulózyalebochronickejhepatitídyB,covyžadujeprísnyskríningpredzahájenímliecby. Zvýšená je aj incidencia herpes zoster. Vakcinácia proti chrípke a pneumokokom je preto kl’úcovoupreventívnoustratégiou(Macejová,2019).


\subsection{Pl’úcneahematologickékomplikácie}

Pl’úcne prejavy sú castou viscerálnou manifestáciou RA. Môžu mat’ formu intersti- ciálnejpl’úcnejchoroby(ILD),ktoráprogredujedonezvratnejfibrózy,pleurálnychvýpotkov aleboreumatoidnýchuzlovvpl’úcnomparenchýme(vkombináciispneumokoniózouznáme ako Caplanov syndróm). Klinickým korulátom je progredujúca dyspnoe a suchý, dráždivý kašel’. V krvnom obraze dominuje normocytárna anémia chronických chorôb, mechaniz- momktorejjehepcidínomsprostredkovanáblokádametabolizmuželeza.Zriedkavou,noži- vot ohrozujúcou entitou je Feltyho syndróm (triáda: RA, splenomegália, granulocytopénia), ktorýextrémnezvyšujenáchylnost’nafulminantnébakteriálneinfekcie(Rovenský,2001). Intersticiálneochoreniepl’úcpriRA Intersticiálne ochorenie pl’úc asociované s reumatoidnou artritídou (RA-ILD) pred- stavujezávažnúextraartikulárnumanifestáciu,ktorávýznamnezhoršujeprognózupacientov. Prevalencia subklinického postihnutia pl’úc detegovaného pomocou vysokorozlišovacieho 10 CT (HRCT) dosahuje až 60 percent, pricom klinicky manifestná forma sa vyskytuje u 10 až 30percentchorých.Rizikovéfaktoryzahrnajúmužsképohlavie,vyššívek,fajcenie,vysoké titreACPAprotilátokaprítomnost’reumatoidnýchuzlov(Kaduraetal.,2021). PatogenézaRA-ILDzahrnakomplexnúinterakciugenetickýchfaktorov,autoimunit- ných mechanizmov a environmentálnych vplyvov. Cytokíny produkované v rámci systémo- vého zápalu, najmä TNF-alfa, IL-6 a transformujúci rastový faktor beta (TGF-β), stimulujú proliferáciu fibroblastov a nadmernú depozíciu kolagénu v pl’úcnom interstíciu. Histopato- logicky prevláda obraz obvyklej intersticiálnej pneumónie (UIP) alebo nešpecifickej inters- ticiálnejpneumónie(NSIP)(Kaduraetal.,2021;Klener,2011). Manažment RA-ILD vyžaduje multidisciplinárny prístup zahrnajúci reumatológa a pneumológa. Terapeutické možnosti sú limitované; niektoré DMARDs (napr. metotrexát) môžu potenciálne zhoršit’ pl’úcne postihnutie, zatial’ co iné (rituximab, abatacept) sa javia ako bezpecnejšie. V pokrocilých štádiách prichádza do úvahy antifibrotická liecba pirfeni- dónomalebonintedanibom(Kaduraetal.,2021).


\subsection{MetabolicképoruchypriRA}

Vzt’ahmedzireumatoidnouartritídouametabolickýmiporuchamijebidirecionálnya komplexný.ChronickýzápalpriRAnegatívneovplyvnujemetabolicképarametre,zatial’co metabolický syndróm zhoršuje zápalovú aktivitu a komplikuje liecbu základného ochorenia (Kašperováetal.,2021). Inzulínovárezistenciaadiabetesmellitus Pacienti s RA majú zvýšené riziko rozvoja diabetes mellitus 2. typu v porovnaní s bežnou populáciou. Prozápalové cytokíny, predovšetkým TNF-α a IL-6, interferujú so sig- nálnou dráhou inzulínu na úrovni inzulínového receptorového substrátu (IRS-1), co vedie k rozvoju inzulínovej rezistencie. Situáciu d’alej komplikuje dlhodobá kortikoterapia, ktorá priamozvyšujeglykémiuapodporujeviscerálnuobezitu(Kašperováetal.,2021). Dyslipidémiaaateroskleróza Lipidový profil u pacientov s aktívnou RA vykazuje charakteristické zmeny ozna- cované ako "lipidový paradox". Napriek nízkym hodnotám celkového cholesterolu a LDL- cholesterolu (ktoré sú konzumované zápalovým procesom) je kardiovaskulárne riziko zvý- šené.Prícinousúkvalitatívnezmenylipoproteínov--oxidáciaLDLcastícazníženáantiatero- génna aktivita HDL cholesterolu. Efektívna kontrola zápalovej aktivity pomocou DMARDs vedie k paradoxnému zvýšeniu hladín lipidov, co však reflektuje zlepšenie metabolického stavu(Freedmanetal.,2014;Kašperováetal.,2021). 11 Reumatickákachexia Osobitným metabolickým fenoménom pri RA je reumatická kachexia, charakteri- zovaná stratou svalovej hmoty pri zachovanej alebo dokonca zvýšenej tukovej hmote. Na rozdielodklasickejkachexieprimalígnychochoreniach,telesnáhmotnost’môžezostat’ne- zmenená, co maskuje závažný úbytok svalstva. Patofyziologicky sa uplatnuje katabolický efekt prozápalových cytokínov, znížená fyzická aktivita pre bolest’ a artritídu a niekedy aj nedostatocný príjem bielkovín. Dôsledkom je svalová slabost’, zhoršená funkcná kapacita a zvýšenérizikopádov(Klener,2011). Dopad reumatoidnej artritídy na kvalitu života (QoL) je devastacný a multodimen- zionálny, zasahujúci d’aleko za hranice fyzickej disability. WHO definuje zdravie ako stav kompletnej bio-psycho-sociálnej pohody, co je ideál, ktorý je u pacientov s RA casto nedo- siahnutel’ný(Krchnaváetal.,2018).


\subsection{Fyzickáafunkcnádoména}

Kl’úcovýmifaktormi,ktorénegatívnedeterminujúQoL,súperzistujúcabolest’,pato- logickáúnavaarannástuhnutost’.Bolest’,prítomnávpokojiajpriaktivite,narúšaspánkovú architektúruavedieksyndrómuchronickéhovycerpania(fatigue),ktorýpacienticastohod- notia ako viac limitujúci než bolest’ samotnú. Ranná stuhnutost’, trvajúca nezriedka hodiny, blokujerannérutinnécinnosti.Progresívnastrataklbovýchrozsahovasvalovejsilyznemož- nujesebaobsluhu(hygiena,obliekanie),covediekzávislostinapomocidruhejosoby,strate autonómieapocitombezmocnosti.Vznikátakcirculusvitiosus:bolest’--inaktivita--atrofia --zhoršeniefunkcie(Krchnaváetal.,2018).


\subsection{Psychologickádoména}

Chronický, celoživotný charakter choroby s nepredvídatel’ným striedaním relapsov a remisií, spolu s hrozbou trvalého postihnutia, generuje u pacientov chronickú úzkost’ a depresívne poruchy. Prevalencia depresie sa odhaduje na 15 -- 40\%, pricom je známe, že depresia znižuje prah bolesti a zhoršuje adherenciu k liecbe. Významným faktorom je aj alterácia telesného obrazu (body image) v dôsledku viditel’ných deformít, co vedie k so- ciálnemustiahnutiu.Pacienticastozažívajúfenomén"naucenejbezmocnostiästratuidentity (Olejárová;kol.,2016).


\subsection{Sociálnaaekonomickádoména}

RA postihuje jedincov na vrchole produktívneho veku, co vedie k vysokej miere in- validizácie a strate pracovnej schopnosti. Absencia príjmu prináša ekonomickú instabilitu. Znížená mobilita vedie k sociálnej izolácii a zanechaniu zál’ub. Ochorenie mení dynamiku rodinných vzt’ahov, ked’ sú partneri nútení prevziat’ rolu opatrovatel’ov. Úlohou fyziotera- pie je prostredníctvom edukácie a tréningu kompenzacných mechanizmov maximalizovat’ funkcnúnezávislost’atýmzlepšit’QoL(Krchnaváetal.,2018;Poliakováetal.,2019). 12


\section{Diferenciálnadiagnostika}

Proces diferenciálnej diagnostiky pri RA je intelektuálnou výzvou, ktorá vyžaduje rigorózne vylúcenie iných nozologických jednotiek s prekrývajúcou sa symptomatológiou. Správna a vcasná diferenciácia je sine qua non pre vol’bu adekvátnej terapie, nakol’ko lie- cebnéprístupysaprijednotlivýchochoreniachdiametrálnelíšia(KarelPavelkaetal.,2018).


\subsection{Osteoartróza(OA)}

Osteoartrózajenajcastejšímïmitátorom"RA,predovšetkýmugeriatrickejpopulácie. Na rozdiel od RA, ktorá je primárne zápalovým ochorením, je OA degeneratívny proces spojenýsopotrebovanímchrupavky. • Charakter bolesti: Pri OA je bolest’ mechanická -- zhoršuje sa námahou a v pokoji ustupuje.PriRAjebolest’zápalová,prítomnáajvkl’ude. • Stuhnutost’:PriOAjekrátka,"štartovacia"(do30minút).PriRAjedlhodobá. • Lokalizácia:OApostihujetypickyDIPklby(Heberdenoveuzly),PIPklby(Bouchar- doveuzly)aklbpalcaruky(rhizartróza).RApostihujeMCPklbyazápästie. • Laboratórium: Pri OA sú zápalové markery (CRP, FW) v norme (Olejárová; kol., 2016).


\subsection{Séronegatívnespondyloartritídy}

Ideoskupinuzápalovýchochorení,prektoréjetypickáneprítomnost’reumatoidného faktoraaväzbanaantigénHLA-B27. 1. Psoriatickáartritída(PsA):Môžebyt’klinickyvel’mipodobnáRA.Typickýmidife- renciacnýmiznakmisúprítomnost’psoriázynakoži(hociartritídamôžepredchádzat’ kožnýmprejavom),postihnutieDIPklbov,daktylitída(párkovitýopuchceléhoprsta)a zmenynanechtoch(jamky,onycholýza).NaRTGvidímesúcasneerózieajnovotvorbu kosti(Macejová,2019). 2. Ankylózujúca spondylitída (Bechterevova choroba): Dominuje postihnutie axiál- neho skeletu -- zápalová bolest’ chrbta a sakroiliitída. Periférna artritída je zvycajne asymetrickáatýkasavel’kýchklbovdolnýchkoncatín(Gravaldietal.,2022). 3. Reaktívna artritída: Vzniká v casovej súvislosti s prekonanou urogenitálnou alebo crevnouinfekciou.KlasickáReiterovatriádazahrnaartritídu,uretritíduakonjunktivi- tídu(Schereretal.,2020).


\subsection{Kryštálovéartropatie(Dna)}

Dnavá artritída je metabolické ochorenie spôsobené depozíciou kryštálov nátrium- urátuvklboch. 13 • Priebeh: Typický je akútny, perakútny dnavý záchvat s extrémne silnou bolest’ou, opuchom a zacervenaním, ktorý vzniká náhle, casto v noci ("podagra"-- postihnutie MTPklbupalcanohy). • Chronickáforma:Tofóznadnamôžeimitovat’RAprítomnost’ouuzlov(tofov)aeró- zií, ale rozlíšenie umožní punkcia synoviálnej tekutiny a dôkaz kryštálov (Rovenský, 2001).


\subsection{Systémovýlupuserythematosus(SLE)}

SLEjemultisystémovéautoimunitnéochorenie. • Klby: Artritída pri SLE býva polyartikulárna a symetrická ako pri RA, ale je typicky neerozívna(Jaccoudovaartropatia)--nenicíkost’,hocispôsobujedeformity. • Ostatnépríznaky:Fotosenzitivita,motýl’ovitýexantémnatvári,postihnutieobliciek (nefritída), serozitídy. Typická je prítomnost’ anti-dsDNA a anti-Sm protilátok (Ro- venskýetal.,2015).


\section{Terapiareumatoidnejartritídy}

Manažment RA prešiel v posledných dvoch dekádach revolucnými zmenami. Sú- casná stratégia je založená na koncepte "Treat to Target" (liecba k ciel’u), ktorého im- peratívom je dosiahnutie klinickej remisie alebo aspon stavu nízkej aktivity choroby v co najkratšom case. Tento prístup si vyžaduje proaktívnu úpravu liecby a úzku spoluprácu re- umatológa,pacientaarehabilitacnéhotímu(Smolenetal.,2018).


\subsection{Farmakologickáliecba}

Farmakoterapia zostáva pilierom liecby a delí sa na symptomatickú, ktorá len tlmí prejavy,achorobumodifikujúcu,ktorámenípriebehochorenia. Symptomatickáliecba(NSAaKortikoidy) Tieto lieky poskytujú pacientovi rýchlu úl’avu, ale neriešia podstatu -- nezastavujú klbovúdeštrukciu. • Nesteroidné antiflogistiká (NSA): Efektívne tlmia bolest’ a rannú stuhnutost’ inhi- bíciou enzýmov cyklooxygenázy. Ich dlhodobé užívanie je však limitované rizikom gastropatie(vredy),kardiovaskulárnychpríhodarenálnehopoškodenia. • Glukokortikoidy: Majú silný a rýchly protizápalový efekt. Využívajú sa najmä ako premost’ujúca liecba (bridging therapy) vo vcasných štádiách alebo pri vzplanutiach. Pre riziko závažných nežiaducich úcinkov (osteoporóza, diabetes, hypertenzia, kata- rakta) je snahou podávat’ co najnižšiu úcinnú dávku co najkratší cas (Olejárová; kol., 2016). 14 Chorobumodifikujúcelieky(DMARDs) DMARDs (Disease Modifying Antirheumatic Drugs) zasahujú priamo do patogene- tickýchmechanizmovaspomal’ujúrádiografickúprogresiu. 1. KonvencnésyntetickéDMARDs(csDMARDs):Metotrexátje"zlatýmštandardomä liekom prvej vol’by. Pôsobí ako antagonista kyseliny listovej. K d’alším liekom patria leflunomid,sulfasalazínaantimalariká(hydroxychlorochín). 2. Biologické DMARDs (bDMARDs): Sú to geneticky inžinierované proteíny, ktoré cieleneblokujúšpecifickécytokínyalebobunky.Súindikovanéprizlyhaníkonvencnej liecby.Patrísem: • anti-TNFliecba(Adalimumab,Infliximab), • blokátoryIL-6(Tocilizumab), • B-bunkovádeplécia(Rituximab), • modulátorykostimulácieT-buniek(Abatacept). 3. Cielené syntetické DMARDs (tsDMARDs): Ide o tzv. malé molekuly (inhibítory JAK-kináz--tofacitinib,baricitinib),ktorésapodávajúperorálneablokujúvnútrobun- kovýprenossignálu.Súrovnakoúcinnéakobiologiká(Schereretal.,2020).


\section{Komplexnárehabilitácia}

Rehabilitácia je neoddelitel’nou súcast’ou terapeutického plánu. Zatial’ co farmako- terapia kontroluje zápal, rehabilitácia rieši funkcné dopady ochorenia. Jej ciel’om je udržat’ rozsah pohybu, svalovú silu, zmiernit’ bolest’ a naucit’ pacienta žit’ s chorobou tak, aby bol conajdlhšiesebestacný(Kolár;kol.,2012).


\subsection{Kinezioterapiaaliecebnátelesnávýchova}

Pohybová liecba (LTV) musí byt’ prísne individuálna a rešpektovat’ aktuálnu ak- tivitu ochorenia (akútne vs. chronické štádium). Nedávne systematické prehl’ady a meta- analýzy potvrdzujú, že pravidelné cvicenie je bezpecné a má signifikantný pozitívny vplyv na funkcnú kapacitu a aeróbnu zdatnost’ bez zhoršenia aktivity ochorenia alebo poškodenia klbov(Zhangetal.,2025;Metsiosetal.,2020). Cvicenie v akútnom štádiu V case exacerbácie (vzplanutia) zápalu je prioritou ochrana klbu(Kolár;kol.,2012). • Polohovanie: Klby sa polohujú vo funkcne výhodných postaveniach na prevenciu flekcnýchkontraktúr. 15 • Pasívne pohyby: Vykonávajú sa opatrne do prahu bolesti na udržanie klbovej vôl’e a výživychrupavky. • Izometrické cvicenie: Svalové kontrakcie bez zmeny dlžky svalu sú kl’úcové na pre- venciurýchlejsvalovejatrofie,ktorásprevádzazápal. (Kolár;kol.,2012;Navrátiletal.,2019) Cvicenie v chronickom štádiu a remisii V tomto období je ciel’om obnova funkcie a kondície(Metsiosetal.,2020). 1. Aeróbny tréning: Chôdza, cyklistika, plávanie. Odporúca sa stredná intenzita (60 -- 70\% maximálnej tepovej frekvencie) 30 minút denne. Aeróbne cvicenie znižuje kar- diovaskulárne riziko, ktoré je u pacientov s RA zvýšené, a redukuje únavu (Metsios etal.,2020). 2. Silový tréning: Cvicenie proti odporu (s pružnými t’ahmi Thera-Band, l’ahkými cin- kami)jenevyhnutnénabojprotireumatickejkachexii(úbytkusvalovejhmoty).Štúdie ukazujú, že dynamický silový tréning je bezpecný a vedie k zlepšeniu svalovej sily a funkcnejschopnosti(Zhangetal.,2025). 3. Cvicenie ruky: Špecifické posilnovacie a mobilizacné cvicenia pre drobné klby ruky (programSARAH)zlepšujúsiluúchopuamanuálnuzrucnost’efektívnejšieakobežná starostlivost’(Williamsetal.,2015). High-intensity interval training Vysokointenzívny intervalový tréning (HIIT) predsta- vuje inovatívny prístup, ktorý bol dlho považovaný za nevhodný pre pacientov so zápalo- vými artropatiami. Nedávne randomizované kontrolované štúdie však priniesli prekvapivé zistenia. Štúdia ExeHeart preukázala, že HIIT v primárnej fyzioterapeutickej starostlivosti je bezpecný a efektívny u pacientov so zápalovou artritídou vrátane RA. Výsledky nazna- cujú zlepšenie aeróbnej kapacity a kardiovaskulárnych parametrov bez zhoršenia zápalovej aktivityochorenia(Nordénetal.,2024). Kl’úcom k úspešnej implementácii HIIT u pacientov s RA je individuálna adaptá- cia protokolu. Intervaly vysokej intenzity by mali byt’ kratšie a odpocinkové fázy dlhšie v porovnaní so zdravou populáciou. Cvicenie by malo prebiehat’ na nebolestivých klboch s preferenciou aktivít s nízkou nárazovou zát’ažou, ako je cyklistika alebo plávanie (Nordén etal.,2024;Metsiosetal.,2020). Podporazmenyživotnéhoštýlu Fyzioterapeutomvedenéintervenciezameranénazmenu správania (behavior change interventions) sa ukazujú ako sl’ubný prístup na zvýšenie dlho- dobej adherencie k pohybovej aktivite. Štúdia PIPPRA preukázala realizovatel’nost’ takejto 16 intervencieupacientovsRA.Kombináciaedukácie,stanoveniaindividuálnychciel’ovapra- videlnéhofollow-upuvediekudržatel’némuzvýšeniufyzickejaktivity(Larkinetal.,2024).


\subsection{Efektívnost’fyzioterapiepriRA}

Systematické prehl’ady a meta-analýzy poskytujú vysokú úroven dôkazov pre úcin- nost’ fyzioterapie v manažmente RA. Najnovší systematický prehl’ad analyzujúci randomi- zované kontrolované štúdie dospel k záveru, že fyzioterapeutická intervencia vedie k sig- nifikantnému zmierneniu bolesti a zlepšeniu funkcných parametrov u pacientov s RA (El Ammareetal.,2023). Špecificky pre rehabilitáciu pri reumatických ochoreniach existujú štandardizované protokoly zahrnajúce kombináciu kinezioterapie, fyzikálnych procedúr a edukácie. Kom- plexný rehabilitacný program by mal byt’ zostavený multidisciplinárnym tímom a prispôso- bený individuálnym potrebám pacienta, fáze ochorenia a prítomnosti komorbidít (Jarošová etal.,2018).


\subsection{Fyzikálnaterapia}

Fyzikálna terapia predstavuje integrálnu súcast’ rehabilitacného programu pri RA. Jej základným princípom je využitie fyzikálnych podnetov -- tepla, chladu, elektrických im- pulzov, svetla a mechanickej energie -- na dosiahnutie terapeutických úcinkov v ciel’ových tkanivách. Tieto procedúry primárne slúžia ako príprava pred kinezioterapiou, pricom ich analgezujúci a myorelaxacný efekt umožnuje pacientovi efektívnejšie cvicit’ (Kraft, 2018; Krížová,2014). Termoterapia Aplikácia tepla (hypertermoterapia) nachádza uplatnenie predovšetkým v chronic- kom štádiu ochorenia. Teplo spôsobuje vazodilatáciu, zlepšuje perfúziu tkanív a uvol’nuje svalový spazmus. Najcastejšie využívané modality zahrnajú parafínové zábaly, ktoré sú ob- zvlášt’ vhodné pre drobné klby rúk, a solux lamp. V akútnom štádiu zápalu je naopak in- dikovaná kryoterapia (chlad), ktorá tlmí zápalovú reakciu a znižuje opuch (Vavrová et al., 2016). Elektroliecba Elektroterapeutické modality majú v liecbe RA dlhodobú tradíciu. TENS (transku- tánna elektrická nervová stimulácia) poskytuje analgéziu prostredníctvom vrátkového me- chanizmutlmeniabolesti.Interferencnéprúdy(IFP)kombinujúvýhodystredofrekvencných a nízkofrekvencných prúdov a sú suverénnou modalitou pre hlboko uložené klby. Podrob- nejšie o IFP pojednáva samostatná kapitola tejto práce (Podebradský et al., 2009; Navrátil etal.,2019). 17 Fototerapiaamechanoterapia Nízkohladinovýlaser(LLLT)aultrazvukpodporujúhojivéprocesyvmäkkýchtkani- vách prostredníctvom fotobiomodulacných a termických úcinkov. Tieto modality sú vhodné pri tendosynovitíde, entezopatii a na podporu hojenia po artikulárnych injekciách (Podeb- radskýetal.,2009).


\subsection{Balneoterapiaakúpel’náliecba}

Významnou súcast’ou komplexnej starostlivosti o pacientov s RA je balneoterapia, ktorámávnašichpodmienkachdlhorocnútradíciu.Ideovyužitieprírodnýchliecivýchvôd, peloidov (bahna) a klimatických faktorov. Podl’a najnovšej meta-analýzy, ktorú publikovala Aribietal.(2025),prinášabalneoterapiasignifikantnézlepšenievredukciibolestiazlepšení kvality života (QoL) u pacientov s reumatickými ochoreniami. Efekt nie je len placebo; ter- málne minerálne vody pôsobia chemicky (cez kožu sa vstrebávajú minerály ako síra, jód) a fyzikálne (teplo, vztlak). Verhagen et al. (2015) v Cochrane systematickom prehl’ade potvr- dzujú,že"kúpel’náliecba"(Spatherapy)mápozitívnyvplyvnafunkcnúkapacitupacientovs RA,pricombenefitypretrvávajú3až6mesiacovpoukonceníliecebnéhopobytu.Špecifický význam majú sírne vody (napr. Piešt’any, Trencianske Teplice). Štúdia autorov Zapolski et al.(2024)poukázalanato,žesírnekúpelenielentlmiazápal,alemajúajpozitívnyvplyvna kardiovaskulárneparametre,cojeuRApacientov(castotrpiacichaterosklerózou)kl’úcové. Mechanizmus úcinku zahrna pravdepodobne aj moduláciu oxidacného stresu. Wa˛tor (2020) uvádzajú, že balneochemické vlastnosti vôd ovplyvnujú redoxný potenciál v tkanivách, cím napomáhajúneutralizáciivol’nýchradikálovvznikajúcichprichronickomzápale.


\subsection{Ergoterapiaaochranaklbov}

Ergoterapia ucí pacienta techniky "Joint Protection"-- ako vykonávat’ bežné cinnosti bez zbytocného pret’ažovania klbov. Zahrna používanie dlahovania na prevenciu deformít, úpravu domáceho prostredia a používanie kompenzacných pomôcok (napr. špeciálne otvá- race,zosilnené rúckypríborov), ktoréšetria drobnéklbyruky aumožnujú pacientovi zostat’ nezávislým(Herbertetal.,2022). Rehabilitacné ošetrovatel’stvo doplna ergoterapeutickú intervenciu o nácvik základ- nýchsebaobslužnýchaktivítupacientovst’ažkýmiporuchamihybnosti.Ciel’omjemaxima- lizácia funkcnej nezávislosti a prevencia sekundárnych komplikácií z imobility (Klusonová etal.,2014).


\subsection{Rehabilitáciaareumatológia--interdisciplinárnaspolupráca}

Optimálny manažment pacienta s RA vyžaduje koordinovanú spoluprácu reumato- lóga, fyzioterapeuta, ergoterapeuta, rehabilitacného lekára a d’alších špecialistov. Rehabi- litacná intervencia by mala byt’ integrovaná do liecebného plánu od momentu stanovenia diagnózyanieažporozvojiireverzibilnýchdeformít(Gúth,2010). 18 Vcasná rehabilitácia zahrna nielen kinezioterapiu a fyzikálnu liecbu, ale aj edukáciu pacienta o ochorení, nácvik kompenzacných stratégií a psychologickú podporu. Pacient by mal byt’ aktívnym partnerom v liecebnom procese a mat’ realistické ocakávania ohl’adom priebehuochoreniaamožnostíterapie(Gúth,2010;Špiritovic,2018).


\section{Prognózareumatoidnejartritídy}

PrognózaRAjevysokovariabilnáazávisíodmnožstvafaktorov.Identifikáciaprog- nostických markerov umožnuje stratifikovat’ pacientov podl’a rizika a prispôsobit’ intenzitu liecby(Wolfeetal.,2010).


\subsection{Faktoryovplyvnujúcedlhodobývýsledok}

Nepriaznivéprognostickéfaktoryzahrnajúvysokétitrereumatoidnéhofaktoraaanti- CCP protilátok, prítomnost’ erozívnych zmien v skorom štádiu ochorenia, extraartikulárne manifestácie,vysokúzápalovúaktivituafunkcnúdisabilitunazaciatkuchoroby.Ženymajú všeobecne horšiu prognózu než muži a genetické faktory (HLA-DR4) sú spojené s agresív- nejšímpriebehom(Wolfeetal.,2010). Vcasné zahájenie liecby DMARDs, dosiahnutie remisie do 6 mesiacov od zaciatku symptómov a striktná kontrola zápalovej aktivity ("Treat to Target") signifikantne zlepšujú dlhodobé výsledky. Pacienti, ktorí dosiahnu remisiu v ökne príležitosti", majú dobré šance na dlhodobé funkcné zachovanie bez významného štrukturálneho poškodenia (K. Pavelka, 2012;McCarthyetal.,2018).


\subsection{Mortalitaakvalitaživota}

